% ============================================================
% Chapter 10: Security Testing: SAST, DAST, Vulnerability Assessment and Penetration Testing
% Language: Greek — edit freely; regenerate from src/content if preferred.
% ============================================================
\chapter{Έλεγχος Ασφάλειας: SAST, DAST, Αξιολόγηση Ευπαθειών και Δοκιμές Διείσδυσης}\label{ch:10}
\chaptermeta{Στατική ανάλυση · Δυναμική ανάλυση · Κύκλος διαχείρισης ευπαθειών · Μεθοδολογία και δεοντολογία δοκιμών διείσδυσης · CVSS και επαγγελματική αναφορά}{Επίπεδο: Μεσαίο–Προχωρημένο επίπεδο · Διασφάλιση}{Χρόνος μελέτης: 7–9 ώρες}

Το Κεφάλαιο 9 περιέγραψε πώς αναπτύσσεται λογισμικό με ασφάλεια· το παρόν κεφάλαιο εξετάζει πώς μπορούμε να \emph{επαληθεύσουμε} ότι είναι ασφαλές. Καμία μεμονωμένη τεχνική δεν εντοπίζει κάθε αδυναμία. Η στατική ανάλυση διαβάζει τον κώδικα χωρίς να τον εκτελεί, η δυναμική ανάλυση εξετάζει μια εφαρμογή σε λειτουργία από έξω, η αξιολόγηση ευπαθειών ελέγχει ολόκληρο ένα περιβάλλον για γνωστές αδυναμίες και η δοκιμή διείσδυσης αποδεικνύει τι θα μπορούσε πράγματι να επιτύχει ένας ικανός αντίπαλος. Καθεμιά απαντά σε διαφορετικό ερώτημα, με διαφορετικό κόστος και σε διαφορετικό στάδιο.

Επειδή ορισμένες από αυτές τις δραστηριότητες χρησιμοποιούν πραγματικά επιθετικές τεχνικές, το κεφάλαιο δίνει ίση βαρύτητα στη μέθοδο και στη δεοντολογία. Εξηγεί τα νομικά θεμέλια του εξουσιοδοτημένου ελέγχου, τα όρια της δραστηριότητας μετά την εκμετάλλευση και τον χειρισμό ευαίσθητων αποδεικτικών στοιχείων, και δείχνει πώς πρέπει να βαθμολογούνται και να αναφέρονται τα ευρήματα, ώστε να οδηγούν σε αποκατάσταση και όχι σε ένα έγγραφο που δεν διαβάζει κανείς.

\begin{outcomesbox}{Μαθησιακά αποτελέσματα}
Μετά τη μελέτη του κεφαλαίου θα πρέπει να μπορείτε:
\begin{itemize}
  \item Να συγκρίνετε τα SAST, DAST, την αξιολόγηση ευπαθειών και τη δοκιμή διείσδυσης ως προς το εύρος, τα πλεονεκτήματα και τους περιορισμούς τους.
  \item Να περιγράφετε τον κύκλο διαχείρισης ευπαθειών και να ιεραρχείτε τα ευρήματα με βάση το CVSS και την πιθανότητα εκμετάλλευσης.
  \item Να εξηγείτε τα μοντέλα, τις φάσεις και τις αναγνωρισμένες μεθοδολογίες των δοκιμών διείσδυσης.
  \item Να διατυπώνετε τις νομικές και δεοντολογικές απαιτήσεις του εξουσιοδοτημένου ελέγχου, συμπεριλαμβανομένων των κανόνων εμπλοκής.
  \item Να συντάσσετε ένα εύρημα που συνδυάζει τεκμηρίωση, αξιολόγηση κινδύνου και εφαρμόσιμη πρόταση αποκατάστασης.
\end{itemize}
\end{outcomesbox}

\section{Στατικός έλεγχος ασφάλειας εφαρμογών (SAST)}\label{sec:10.1}
Ο \textbf{στατικός έλεγχος ασφάλειας εφαρμογών} (Static Application Security Testing, SAST) αναλύει πηγαίο κώδικα, ενδιάμεσο κώδικα ή εκτελέσιμα \emph{χωρίς να τα εκτελεί}. Είναι τεχνική «λευκού κουτιού»: το εργαλείο έχει πρόσβαση σε ολόκληρο τον κώδικα. Τα σύγχρονα εργαλεία SAST κατασκευάζουν ένα μοντέλο του προγράμματος και εκτελούν \textbf{ανάλυση μόλυνσης} (taint analysis), παρακολουθώντας τα δεδομένα από τις \emph{πηγές} (όπως οι παράμετροι HTTP) έως τα επικίνδυνα \emph{σημεία κατάληξης} (όπως η εκτέλεση SQL ή η έξοδος HTML). Αν μολυσμένα δεδομένα φτάσουν σε σημείο κατάληξης χωρίς να περάσουν από αναγνωρισμένο \emph{μηχανισμό εξυγίανσης}, το εργαλείο αναφέρει πιθανή ευπάθεια μαζί με την ακριβή γραμμή του κώδικα.

Το μεγάλο πλεονέκτημα του SAST είναι ο χρόνος εφαρμογής του: μπορεί να εκτελείται στον επεξεργαστή κώδικα του προγραμματιστή και σε κάθε υποβολή, πολύ πριν εγκατασταθεί η εφαρμογή. Η κύρια αδυναμία του είναι η ακρίβεια. Επειδή δεν μπορεί να παρατηρήσει τη συμπεριφορά κατά την εκτέλεση, παράγει \textbf{ψευδώς θετικά} αποτελέσματα (αναφερόμενα ζητήματα που δεν είναι εκμεταλλεύσιμα) και παραλείπει σφάλματα που εξαρτώνται από τη διαμόρφωση ή από την επιχειρησιακή λογική. Τα επιτυχημένα προγράμματα προσαρμόζουν επομένως τους κανόνες στα πλαίσια ανάπτυξης που χρησιμοποιούν, αξιολογούν τα αποτελέσματα και αντιμετωπίζουν το SAST ως ένα από πολλά επίπεδα.

\section{Δυναμικός έλεγχος ασφάλειας εφαρμογών (DAST)}\label{sec:10.2}
Ο \textbf{δυναμικός έλεγχος ασφάλειας εφαρμογών} (Dynamic Application Security Testing, DAST) εξετάζει μια εφαρμογή \emph{σε λειτουργία} από έξω, όπως θα έκανε ένας επιτιθέμενος, χωρίς πρόσβαση στον πηγαίο κώδικα. Ένας σαρωτής DAST, όπως το OWASP ZAP, πρώτα \emph{ανιχνεύει} την εφαρμογή για να ανακαλύψει σελίδες, παραμέτρους και σημεία πρόσβασης API και στη συνέχεια αποστέλλει ειδικά διαμορφωμένες εισόδους και αναλύει τις αποκρίσεις για ενδείξεις ευπάθειας —ένα μήνυμα σφάλματος SQL, ένα ανακλώμενο σενάριο, μια κεφαλίδα ασφάλειας που λείπει. Επειδή παρατηρεί την πραγματική συμπεριφορά, ένα εύρημα DAST είναι συνήθως πράγματι εκμεταλλεύσιμο, και μπορεί να εντοπίσει προβλήματα διαμόρφωσης και διακομιστή που το SAST δεν βλέπει.

Το DAST έχει επίσης όρια. Μπορεί να ελέγξει μόνο όσα καταφέρνει να προσεγγίσει, οπότε σελίδες πίσω από σύνθετη ταυτοποίηση ή ροές πολλών βημάτων ενδέχεται να παραλειφθούν, και δεν μπορεί να υποδείξει την υπεύθυνη γραμμή κώδικα. Εκτελείται σε όψιμο στάδιο του κύκλου ζωής, σε εγκατεστημένο περιβάλλον δοκιμών. Ο \textbf{διαδραστικός έλεγχος} (Interactive Application Security Testing, IAST) συνδυάζει τις δύο προσεγγίσεις, εξοπλίζοντας την εφαρμογή με αισθητήρες ενώ εκτελούνται οι δυναμικές δοκιμές, ώστε κάθε παρατηρούμενο ζήτημα να συνδέεται με τη θέση του στον κώδικα.

\section{Ο κύκλος διαχείρισης ευπαθειών}\label{sec:10.3}
Η \textbf{αξιολόγηση ευπαθειών} εντοπίζει συστηματικά γνωστές αδυναμίες —ενημερώσεις που λείπουν, μη ασφαλείς διαμορφώσεις, εργοστασιακά διαπιστευτήρια— σε δίκτυα, υπολογιστές και εφαρμογές, συνήθως με αυτοματοποιημένους σαρωτές όπως τα Nessus, OpenVAS ή τα σενάρια του Nmap. Είναι εκτενής αλλά ρηχή: αναφέρει τι \emph{ενδέχεται} να είναι εκμεταλλεύσιμο χωρίς να επιχειρεί εκμετάλλευση. Η αξιολόγηση αποτελεί ένα στάδιο του συνεχούς \textbf{κύκλου διαχείρισης ευπαθειών}: \emph{ανακάλυψη περιουσιακών στοιχείων} (δεν μπορείς να προστατεύσεις ό,τι δεν γνωρίζεις ότι έχεις), \emph{σάρωση}, \emph{ανάλυση και ιεράρχηση}, \emph{αποκατάσταση}, \emph{επαλήθευση} ότι η διόρθωση λειτουργεί και \emph{αναφορά} των τάσεων.

Η ιεράρχηση είναι το σημείο όπου δυσκολεύονται τα περισσότερα προγράμματα, επειδή οι σαρωτές αναφέρουν συνήθως χιλιάδες ευρήματα. Η σοβαρότητα από μόνη της δεν αρκεί. Μια τεκμηριωμένη προσέγγιση συνδυάζει την τεχνική σοβαρότητα (CVSS), την πιθανότητα εκμετάλλευσης —για παράδειγμα τη βαθμολογία \textbf{EPSS} ή την καταχώριση στον κατάλογο \emph{Known Exploited Vulnerabilities} της CISA— και την επιχειρησιακή σημασία και έκθεση του επηρεαζόμενου στοιχείου. Ένας διακομιστής εκτεθειμένος στο διαδίκτυο με ευπάθεια που εκμεταλλεύονται ενεργά οι επιτιθέμενοι χρειάζεται προσοχή σήμερα, ακόμη κι αν ένα «κρίσιμο» εύρημα σε ένα απομονωμένο μηχάνημα δοκιμών περιμένει.

\section{Δοκιμές διείσδυσης: μοντέλα, φάσεις, δεοντολογία και πλαίσια}\label{sec:10.4}
Μια \textbf{δοκιμή διείσδυσης} είναι μια εξουσιοδοτημένη, προσομοιωμένη επίθεση που επιχειρεί να εκμεταλλευτεί ευπάθειες, ώστε να αποδείξει την πραγματική τους επίπτωση. Οι δοκιμές κατατάσσονται με βάση τη γνώση που δίνεται στον ελεγκτή: \textbf{μαύρου κουτιού} (χωρίς προηγούμενη πληροφόρηση, όπως ένας εξωτερικός επιτιθέμενος), \textbf{γκρίζου κουτιού} (ορισμένες πληροφορίες, όπως ένας λογαριασμός χρήστη) και \textbf{λευκού κουτιού} (πλήρης τεκμηρίωση και πηγαίος κώδικας, που προσφέρει την πληρέστερη κάλυψη στον διαθέσιμο χρόνο). Μια τυπική ανάθεση εξελίσσεται σε φάσεις: \emph{προετοιμασία} και καθορισμός πεδίου, \emph{αναγνώριση}, \emph{σάρωση και απαρίθμηση}, \emph{εκμετάλλευση}, \emph{δραστηριότητα μετά την εκμετάλλευση} και \emph{αναφορά}. Αναγνωρισμένες μεθοδολογίες, όπως οι PTES, OWASP Web Security Testing Guide, OSSTMM και NIST SP 800-115, δομούν αυτές τις φάσεις και καθιστούν τα αποτελέσματα επαναλήψιμα.

Η νομιμότητα στηρίζεται αποκλειστικά στην \textbf{εξουσιοδότηση}. Πριν ξεκινήσει οποιαδήποτε δοκιμή, ο πελάτης πρέπει να υπογράψει σύμβαση και \textbf{κανόνες εμπλοκής} (rules of engagement) που ορίζουν το πεδίο (ποια συστήματα, διευθύνσεις και εφαρμογές), τις επιτρεπόμενες τεχνικές, τα χρονικά παράθυρα, τα πρόσωπα επικοινωνίας έκτακτης ανάγκης και τον τρόπο χειρισμού ευαίσθητων δεδομένων. Ο έλεγχος ενός συστήματος εκτός του συμφωνημένου πεδίου —ακόμη κι ενός που φαίνεται προφανώς σχετικό— αποτελεί μη εξουσιοδοτημένη πρόσβαση. Οι επαγγελματίες ελεγκτές αποφεύγουν επίσης τις καταστροφικές ενέργειες και σταματούν αμέσως αν εντοπίσουν ενδείξεις πραγματικής, εν εξελίξει παραβίασης.

\section{Όρια της φάσης μετά την εκμετάλλευση και χειρισμός αποδεικτικών στοιχείων}\label{sec:10.5}
Η φάση \textbf{μετά την εκμετάλλευση} (post-exploitation) αποδεικνύει την επιχειρησιακή επίπτωση μιας παραβίασης: ποια δεδομένα θα μπορούσαν να προσπελαστούν, αν ήταν δυνατή η κλιμάκωση προνομίων και πόσο μακριά θα μπορούσε να κινηθεί ο ελεγκτής μέσα στο δίκτυο. Είναι επίσης η φάση με τον μεγαλύτερο δεοντολογικό κίνδυνο. Ο κατευθυντήριος κανόνας είναι η \textbf{ελάχιστη αναγκαία ενέργεια}: αποδεικνύεις την πρόσβαση, δεν την εκμεταλλεύεσαι. Ένας ελεγκτής που φτάνει σε μια βάση δεδομένων πελατών καταγράφει ένα στιγμιότυπο της δομής του πίνακα και τον αριθμό των εγγραφών, όχι τις ίδιες τις εγγραφές. Κάθε μηχανισμός μόνιμης παρουσίας, λογαριασμός ή αρχείο που δημιουργείται κατά τη δοκιμή πρέπει να τεκμηριώνεται και να αφαιρείται στη συνέχεια.

Τα αποδεικτικά στοιχεία που συλλέγονται κατά τη δοκιμή —στιγμιότυπα, υποκλαπέντα διαπιστευτήρια, αποσπάσματα ευαίσθητων δεδομένων— είναι και τα ίδια ευαίσθητα. Πρέπει να αποθηκεύονται κρυπτογραφημένα, να κοινοποιούνται μόνο σε εξουσιοδοτημένους παραλήπτες, να διατηρούνται μόνο για όσο διάστημα απαιτεί η σύμβαση και στη συνέχεια να καταστρέφονται με ασφάλεια. Όταν εμπλέκονται προσωπικά δεδομένα, η επεξεργασία πρέπει επιπλέον να συμμορφώνεται με τη νομοθεσία προστασίας δεδομένων, όπως ο ΓΚΠΔ (Κεφάλαιο 13).

\section{Βαθμολόγηση και αναφορά: CVSS, τεκμηρίωση και αποκατάσταση}\label{sec:10.6}
Το \textbf{Common Vulnerability Scoring System (CVSS)} εκφράζει την τεχνική σοβαρότητα μιας ευπάθειας με βαθμολογία από 0 έως 10. Οι \emph{βασικές μετρικές} του περιγράφουν πώς μπορεί να γίνει εκμετάλλευση της ευπάθειας (διάνυσμα επίθεσης, πολυπλοκότητα, απαιτούμενα προνόμια, αλληλεπίδραση χρήστη) και ποια είναι η επίπτωσή της στην εμπιστευτικότητα, την ακεραιότητα και τη διαθεσιμότητα. Η βαθμολογία συνοδεύεται από μια \emph{συμβολοσειρά διανύσματος}, όπως \texttt{CVSS:3.1/AV:N/AC:L/PR:N/UI:N/S:U/C:H/I:H/A:H} (βαθμολογία 9,8), η οποία καθιστά διαφανή τη συλλογιστική. Το CVSS μετρά τη σοβαρότητα και όχι τον κίνδυνο· το πλαίσιο του οργανισμού πρέπει πάντοτε να λαμβάνεται υπόψη.

Μια επαγγελματική αναφορά απευθύνεται σε δύο κοινά. Η \textbf{σύνοψη για τη διοίκηση} εξηγεί, σε μη τεχνική γλώσσα, τον συνολικό κίνδυνο και τις σημαντικότερες συστάσεις. Κάθε \textbf{τεχνικό εύρημα} περιλαμβάνει σαφή τίτλο, τα επηρεαζόμενα στοιχεία, περιγραφή, τεκμηρίωση \emph{αναπαραγωγής} βήμα προς βήμα, τη σοβαρότητα με το διάνυσμα CVSS, την επιχειρησιακή επίπτωση και μια συγκεκριμένη, εφαρμόσιμη \textbf{πρόταση αποκατάστασης} —όχι «βελτιώστε την ασφάλεια», αλλά «αντικαταστήστε τη συνένωση συμβολοσειρών στη γραμμή 42 του \texttt{search.php} με παραμετροποιημένο ερώτημα και εφαρμόστε έναν κανόνα WAF ως προσωρινό μέτρο». Ένας επανέλεγχος επαληθεύει στη συνέχεια ότι οι διορθώσεις είναι αποτελεσματικές.

\begin{table}[htbp]
  \centering\small
  \caption{Σύγκριση προσεγγίσεων ελέγχου}
  \begin{tabularx}{\linewidth}{>{\raggedright\arraybackslash}X >{\raggedright\arraybackslash}X >{\raggedright\arraybackslash}X >{\raggedright\arraybackslash}X >{\raggedright\arraybackslash}X}
    \toprule
    \textbf{Προσέγγιση} & \textbf{Πρόσβαση} & \textbf{Πότε} & \textbf{Πλεονέκτημα} & \textbf{Περιορισμός} \\
    \midrule
    SAST & Πηγαίος κώδικας (λευκό κουτί) & Κατά την ανάπτυξη & Πρώιμο, υποδεικνύει την ακριβή γραμμή & Ψευδώς θετικά, καμία εικόνα εκτέλεσης \\
    DAST & Εφαρμογή σε λειτουργία (μαύρο κουτί) & Περιβάλλον δοκιμών & Επιβεβαιώνει την πραγματική συμπεριφορά & Περιορισμένη κάλυψη, όψιμο \\
    Αξιολόγηση ευπαθειών & Δίκτυο, υπολογιστές, εφαρμογές & Συνεχώς & Ευρεία κάλυψη & Καμία απόδειξη εκμεταλλευσιμότητας \\
    Δοκιμή διείσδυσης & Συμφωνημένο πεδίο & Περιοδικά / μετά από σημαντική αλλαγή & Αποδεικνύει την πραγματική επίπτωση & Χρονικά περιορισμένη, δαπανηρή, στιγμιαία εικόνα \\
    \bottomrule
  \end{tabularx}
\end{table}
\section*{Βασικοί όροι}
\begin{description}[style=nextline,leftmargin=1.2em]
  \item[Ανάλυση μόλυνσης] Παρακολούθηση μη αξιόπιστων δεδομένων από τις πηγές έως τα επικίνδυνα σημεία κατάληξης στον κώδικα.
  \item[Ψευδώς θετικό] Αναφερόμενο ζήτημα που στην πραγματικότητα δεν είναι εκμεταλλεύσιμο.
  \item[EPSS] Σύστημα πρόβλεψης εκμετάλλευσης: η εκτιμώμενη πιθανότητα να γίνει εκμετάλλευση μιας ευπάθειας.
  \item[Κανόνες εμπλοκής] Υπογεγραμμένο έγγραφο που ορίζει το πεδίο, τις μεθόδους, τον χρόνο και τα πρόσωπα επικοινωνίας μιας δοκιμής.
  \item[Μετά την εκμετάλλευση] Δραστηριότητες μετά την αρχική πρόσβαση που αποδεικνύουν την επίπτωση, με όριο την ελάχιστη αναγκαία ενέργεια.
  \item[Διάνυσμα CVSS] Συμβολοσειρά που καταγράφει κάθε μετρική που χρησιμοποιήθηκε για τον υπολογισμό μιας βαθμολογίας CVSS.
\end{description}

\section*{Σύνοψη κεφαλαίου}
\begin{itemize}
  \item Τα SAST, DAST, η αξιολόγηση ευπαθειών και η δοκιμή διείσδυσης απαντούν σε διαφορετικά ερωτήματα και αλληλοσυμπληρώνονται.
  \item Η διαχείριση ευπαθειών είναι συνεχής κύκλος· η ιεράρχηση πρέπει να συνδυάζει σοβαρότητα, πιθανότητα εκμετάλλευσης και πλαίσιο του στοιχείου.
  \item Οι δοκιμές διείσδυσης ακολουθούν αναγνωρισμένες φάσεις και μεθοδολογίες και είναι νόμιμες μόνο εντός υπογεγραμμένου πεδίου και κανόνων εμπλοκής.
  \item Η φάση μετά την εκμετάλλευση πρέπει να αποδεικνύει την πρόσβαση με την ελάχιστη αναγκαία ενέργεια, ενώ τα αποδεικτικά στοιχεία της δοκιμής προστατεύονται όπως τα δεδομένα παραγωγής.
  \item Οι καλές αναφορές συνδυάζουν διαφανή βαθμολόγηση CVSS με αναπαραγώγιμη τεκμηρίωση και συγκεκριμένη αποκατάσταση.
\end{itemize}

\section*{Ερωτήσεις ανασκόπησης}
\begin{enumerate}
  \item Γιατί το SAST μπορεί να αναφέρει μια ευπάθεια που το DAST δεν μπορεί να επιβεβαιώσει, και αντιστρόφως;
  \item Δύο ευρήματα: CVSS 9,1 σε έναν απομονωμένο εργαστηριακό διακομιστή και CVSS 7,5 σε έναν διακομιστή εκτεθειμένο στο διαδίκτυο που περιλαμβάνεται στον κατάλογο CISA KEV. Ποιο διορθώνετε πρώτο και γιατί;
  \item Κατά τη διάρκεια μιας δοκιμής ανακαλύπτετε ότι ένα σύστημα λίγο έξω από το πεδίο σας είναι ευάλωτο. Τι πρέπει να κάνετε;
  \item Αναδιατυπώστε τη σύσταση «Διορθώστε το πρόβλημα XSS» ώστε να γίνει εφαρμόσιμη πρόταση αποκατάστασης.
\end{enumerate}
